% GENERATED FILE. Edit canonical lessons, then run npm run generate:books.
% canonical-source-hash: fnv1a64:7ee0f6418b407312
% canonical-lessons: TE-C325-adharapadu, TE-C325-adharapadi-untundi, TE-C325-undi-unte, TE-W325-six-lines, TE-R325-depends-and-would

\chapter{It Depends, and What Would Have Been}
\label{ch:te-it-depends-and-what-would-have-been}
\hlchaptermodality{325}

\begin{chapteropening}
\textbf{By the end of this chapter:} \emph{I can say what something depends on, say what would have happened with -i uṇṭē … -ēvāḍini, and write six connected lines about a trip.}

Complete TE-R325-depends-and-would: say what something depends on and what would have been, after writing six connected lines about a trip.
\end{chapteropening}

\section[\texorpdfstring{\te{ఆధారపడు} (ādhārapaḍu)}{ఆధారపడు (ādhārapaḍu)}]{\te{ఆధారపడు} (ādhārapaḍu) — to depend, to rely}

\label{lesson:TE-C325-adharapadu}



\begin{quote}
Before the new one: say \textquotedblleft{}a support, a basis\textquotedblright{} (\emph{ādhāraṁ}) and \textquotedblleft{}to hurry\textquotedblright{} (\emph{tondarapaḍu}). Then \textquotedblleft{}if it rains tomorrow, I will not come\textquotedblright{}. (\emph{rēpu varṣaṁ paḍitē nēnu rānu}.)
\end{quote}

\subsection*{You'll want to know: \te{ఆధారపడు}}
\textbf{\te{ఆధారపడు}} — \emph{ādhārapaḍu} — \textquotedblleft{}to depend (on), to rely (on)\textquotedblright{}.

\begin{quote}
\textbf{\te{అమ్మమ్మ} \te{నాన్న} \te{మీద} \te{ఆధారపడుతుంది}.}
\end{quote}

\begin{quote}
\emph{ammamma nānna mīda ādhārapaḍutundi.} — \textquotedblleft{}Grandmother depends on Father.\textquotedblright{}
\end{quote}

\begin{quote}
\textbf{\te{మా} \te{ఊరు} \te{వర్షం} \te{మీద} \te{ఆధారపడుతుంది}.}
\end{quote}

\begin{quote}
\emph{mā ūru varṣaṁ mīda ādhārapaḍutundi.} — \textquotedblleft{}Our town depends on the rain.\textquotedblright{}
\end{quote}

What you depend on takes \textbf{\te{మీద}} (\emph{mīda}, \textquotedblleft{}on\textquotedblright{}). The past, \textquotedblleft{}I depended\textquotedblright{}, is \textbf{\te{ఆధారపడ్డాను}} (\emph{ādhārapaḍḍānu}): \textbf{-\te{పడు}} makes its past \textbf{\te{పడ్డాను}}, as \textbf{\te{తొందరపడు}} makes \textbf{\te{తొందరపడ్డాను}}.

\begin{cousinweb}
\textbf{\te{ఆధారపడు}} (\emph{ādhārapaḍu}) is two pieces. \textbf{\te{ఆధారం}}, \textquotedblleft{}a support, a basis\textquotedblright{}, which you know, is borrowed from Sanskrit \emph{ādhāra} (CDIAL 1165). \textbf{-\te{పడు}} is the native verb that DEDR 3853 lists as \textquotedblleft{}to feel, to suffer\textquotedblright{}, and that makes verbs of nouns, as in \textbf{\te{తొందరపడు}} and \textbf{\te{సిగ్గుపడు}}. So: \textquotedblleft{}to come to rest on a support\textquotedblright{}.
\end{cousinweb}

\subsection*{Your turn}
Say these aloud:

\begin{itemize}
\raggedright
  \item \emph{ādhārapaḍu}
  \item \textquotedblleft{}Grandmother depends on Father\textquotedblright{} — \emph{ammamma nānna mīda ādhārapaḍutundi}
  \item \textquotedblleft{}Our town depends on the rain\textquotedblright{} — \emph{mā ūru varṣaṁ mīda ādhārapaḍutundi}
  \item \textquotedblleft{}I depended\textquotedblright{} — \emph{ādhārapaḍḍānu}
\end{itemize}

\begin{tcolorbox}[breakable,colback=teal!4,colframe=teal!35!black,title={Before you move on}]
What does \textbf{\te{ఆధారపడు}} mean? (\textbf{To depend, to rely.}) What does the thing you depend on take? (\textbf{\te{మీద}}, \textquotedblleft{}on\textquotedblright{}.) Which noun is inside it? (\textbf{\te{ఆధారం}}, \textquotedblleft{}a support\textquotedblright{}.)
\end{tcolorbox}
\section[\texorpdfstring{\te{మీద} \te{ఆధారపడి} \te{ఉంటుంది} (mīda ādhārapaḍi …)}{మీద ఆధారపడి ఉంటుంది (mīda ādhārapaḍi …)}]{\te{మీద} \te{ఆధారపడి} \te{ఉంటుంది} (mīda ādhārapaḍi uṇṭundi) — it depends}

\label{lesson:TE-C325-adharapadi-untundi}



\begin{quote}
Say \textquotedblleft{}our town depends on the rain\textquotedblright{} (\emph{mā ūru varṣaṁ mīda ādhārapaḍutundi}). Then \textquotedblleft{}according to\textquotedblright{}. (\emph{prakāraṁ}.)
\end{quote}

\subsection*{You'll want to know: … \te{మీద} \te{ఆధారపడి} \te{ఉంటుంది}}
\begin{quote}
\textbf{\te{అది} \te{వాతావరణం} \te{మీద} \te{ఆధారపడి} \te{ఉంటుంది}.}
\end{quote}

\begin{quote}
\emph{adi vātāvaraṇaṁ mīda ādhārapaḍi uṇṭundi.} — \textquotedblleft{}That depends on the weather.\textquotedblright{}
\end{quote}

\begin{quote}
\textbf{\te{రేపటి} \te{ప్రయాణం} \te{వర్షం} \te{మీద} \te{ఆధారపడి} \te{ఉంటుంది}.}
\end{quote}

\begin{quote}
\emph{rēpaṭi prayāṇaṁ varṣaṁ mīda ādhārapaḍi uṇṭundi.} — \textquotedblleft{}Tomorrow's trip depends on the rain.\textquotedblright{}
\end{quote}

\textbf{\te{ఆధారపడి}}, the \textbf{-\te{ఇ}} form, with \textbf{\te{ఉంటుంది}}, \textquotedblleft{}it stays\textquotedblright{}: \textquotedblleft{}it stays resting on\textquotedblright{}. This is the answer to a question that has no single answer: > \textbf{\te{అది} \te{మీ} \te{మీద} \te{ఆధారపడి} \te{ఉంటుంది}.}

\begin{quote}
\emph{adi mī mīda ādhārapaḍi uṇṭundi.} — \textquotedblleft{}That is up to you.\textquotedblright{}
\end{quote}

\textbf{\te{రేపటి}} (\emph{rēpaṭi}, \textquotedblleft{}tomorrow's\textquotedblright{}) is \textbf{\te{రేపు}} before a noun. When the answer comes from a rule or a list, use \textbf{\te{ప్రకారం}}, \textquotedblleft{}according to\textquotedblright{}:

\begin{quote}
\textbf{\te{కాలపట్టిక} \te{ప్రకారం} \te{రైలు} \te{ఆరు} \te{గంటలకు} \te{వస్తుంది}.}
\end{quote}

\begin{quote}
\emph{kālapaṭṭika prakāraṁ railu āru gaṇṭalaku vastundi.} — \textquotedblleft{}According to the timetable, the train comes at six.\textquotedblright{}
\end{quote}

\subsection*{Your turn}
Say these aloud:

\begin{itemize}
\raggedright
  \item \textquotedblleft{}That depends on the weather\textquotedblright{} — \emph{adi vātāvaraṇaṁ mīda ādhārapaḍi uṇṭundi}
  \item \textquotedblleft{}Tomorrow's trip depends on the rain\textquotedblright{} — \emph{rēpaṭi prayāṇaṁ varṣaṁ mīda ādhārapaḍi uṇṭundi}
  \item \textquotedblleft{}That is up to you\textquotedblright{} — \emph{adi mī mīda ādhārapaḍi uṇṭundi}
  \item \textquotedblleft{}According to the timetable, the train comes at six\textquotedblright{} — \emph{kālapaṭṭika prakāraṁ railu āru gaṇṭalaku vastundi}
  \item 'will you come tomorrow?' — answer: it depends, and say on what, with an if
\end{itemize}

\begin{tcolorbox}[breakable,colback=teal!4,colframe=teal!35!black,title={Before you move on}]
\textquotedblleft{}It depends on the weather\textquotedblright{}? (\textbf{\te{వాతావరణం} \te{మీద} \te{ఆధారపడి} \te{ఉంటుంది}}.) \textquotedblleft{}Up to you\textquotedblright{}? (\textbf{\te{మీ} \te{మీద} \te{ఆధారపడి} \te{ఉంటుంది}}.) \textquotedblleft{}According to\textquotedblright{}? (\textbf{\te{ప్రకారం}}.)
\end{tcolorbox}
\section[\texorpdfstring{\te{ఉండి} \te{ఉంటే} (uṇḍi uṇṭē)}{ఉండి ఉంటే (uṇḍi uṇṭē)}]{\te{ఉండి} \te{ఉంటే} (uṇḍi uṇṭē) — what would have been}

\label{lesson:TE-C325-undi-unte}



\begin{quote}
Say \textquotedblleft{}I used to go\textquotedblright{}, as a man and as a woman (\emph{veḷḷēvāḍini}, \emph{veḷḷēdānni}). Then \textquotedblleft{}if there is\textquotedblright{} (\emph{uṇṭē}) and \textquotedblleft{}having come\textquotedblright{} (\emph{vacci}).
\end{quote}

\subsection*{You'll want to know: -\te{ఇ} \te{ఉంటే}, 'if … had'}
\begin{quote}
\textbf{\te{నా} \te{దగ్గర} \te{డబ్బు} \te{ఉండి} \te{ఉంటే}, \te{నేను} \te{వెళ్ళేవాడిని}.}
\end{quote}

\begin{quote}
\emph{nā daggara ḍabbu uṇḍi uṇṭē, nēnu veḷḷēvāḍini.} — \textquotedblleft{}If I had had money, I would have gone.\textquotedblright{}
\end{quote}

Two pieces from this core, joined: the \textbf{-\te{ఇ}} form and \textbf{\te{ఉంటే}}. \textbf{\te{ఉండి} \te{ఉంటే}}, \textquotedblleft{}if there had been\textquotedblright{}; \textbf{\te{వచ్చి} \te{ఉంటే}}, \textquotedblleft{}if (you) had come\textquotedblright{}; \textbf{\te{చెప్పి} \te{ఉంటే}}, \textquotedblleft{}if (you) had told\textquotedblright{}. It looks at something that did not happen. \textbf{\te{నా} \te{దగ్గర}}, \textquotedblleft{}near me\textquotedblright{}, is how Telugu says \textquotedblleft{}I have (it on me)\textquotedblright{}.

\subsection*{You'll want to know: -\te{ఏవాడిని}, 'would have'}
The main verb is the form of the habit, from the first chapter of this core: \textbf{-\te{ఏవాడిని}}, \textbf{-\te{ఏదాన్ని}}. The same form says \textquotedblleft{}used to\textquotedblright{} and \textquotedblleft{}would have\textquotedblright{}; the \textbf{-\te{ఇ} \te{ఉంటే}} clause in front tells you which.

\begin{quote}
\textbf{\te{నువ్వు} \te{ముందే} \te{చెప్పి} \te{ఉంటే}, \te{నేను} \te{వచ్చేదాన్ని}.}
\end{quote}

\begin{quote}
\emph{nuvvu mundē ceppi uṇṭē, nēnu vaccēdānni.} — \textquotedblleft{}If you had told me earlier, I would have come.\textquotedblright{}
\end{quote}

\begin{quote}
\textbf{\te{బస్సు} \te{సమయానికి} \te{వచ్చి} \te{ఉంటే}, \te{మేము} \te{రైలు} \te{అందుకునేవాళ్ళం}.}
\end{quote}

\begin{quote}
\emph{bassu samayāniki vacci uṇṭē, mēmu railu andukunēvāḷḷaṁ.} — \textquotedblleft{}If the bus had come on time, we would have caught the train.\textquotedblright{}
\end{quote}

\textbf{\te{ముందే}} (\emph{mundē}, \textquotedblleft{}earlier, beforehand\textquotedblright{}) is \textbf{\te{ముందు}} with \textbf{-\te{ఏ}}; \textbf{\te{సమయానికి}} (\emph{samayāniki}, \textquotedblleft{}on time\textquotedblright{}) is \textbf{\te{సమయం}} with \textbf{-\te{కి}}.

\subsection*{Your turn}
Say these aloud:

\begin{itemize}
\raggedright
  \item \textquotedblleft{}If I had had money, I would have gone\textquotedblright{} — \emph{nā daggara ḍabbu uṇḍi uṇṭē, nēnu veḷḷēvāḍini}
  \item \textquotedblleft{}A woman says: if you had told me earlier, I would have come\textquotedblright{} — \emph{nuvvu mundē ceppi uṇṭē, nēnu vaccēdānni}
  \item \textquotedblleft{}If the bus had come on time, we would have caught the train\textquotedblright{} — \emph{bassu samayāniki vacci uṇṭē, mēmu railu andukunēvāḷḷaṁ}
  \item one thing from your last trip that would have gone better, and why
\end{itemize}

\begin{tcolorbox}[breakable,colback=teal!4,colframe=teal!35!black,title={Before you move on}]
\textquotedblleft{}If (you) had come\textquotedblright{}? (\textbf{\te{వచ్చి} \te{ఉంటే}}.) \textquotedblleft{}I would have gone\textquotedblright{}? (\textbf{\te{వెళ్ళేవాడిని}}, \textbf{\te{వెళ్ళేదాన్ని}}.) What else does that form say? (\textbf{\textquotedblleft{}Used to\textquotedblright{}.})
\end{tcolorbox}
\section[(six lines)]{Six lines about a trip}

\label{lesson:TE-W325-six-lines}



\begin{quote}
The last A2 writing task asked for a connected message about a familiar day. This one asks for six lines, and they tell a story: a trip, to your home town or anywhere else.

Pick one reader: a friend, by name.
\end{quote}

\subsection*{What to know first}
A told story is not a list of events. It has a \textbf{scene}, in \textbf{-\te{తున్నప్పుడు}} or \textbf{\te{ఉండేది}}. It has \textbf{events} in the plain past, kept in order with \textbf{\te{మొదట}}, \textbf{\te{ఆ} \te{తరువాత}}, \textbf{-\te{ఇ}} and \textbf{\te{చివరికి}}. It has at least one \textbf{reason}, with \textbf{\te{కాబట్టి}} or \textbf{\te{ఎందుకంటే}}, and a \textbf{feeling}, with \textbf{\te{నాకు} … -\te{గా} \te{అనిపించింది}}.

A good one ends by looking at what might have been: \textbf{… -\te{ఇ} \te{ఉంటే}, … -\te{ఏవాడిని}} (or \textbf{-\te{ఏదాన్ని}}).

\subsection*{Writing — connected composition}
On paper, in Telugu script, write \textbf{six lines} to your reader about a trip, real or invented, one line for each of these:

\begin{enumerate}
\raggedright
  \item the scene: when it was, and what you were doing as it began, with \textbf{-\te{తున్నప్పుడు}};
  \item the first thing that happened, with \textbf{\te{మొదట}}, and a ticket or a seat;
  \item what happened next, with \textbf{\te{ఆ} \te{తరువాత}} or an \textbf{-\te{ఇ}} form, and one thing that went wrong;
  \item why, with \textbf{\te{కాబట్టి}} or \textbf{\te{ఎందుకంటే}}; and if you like, what someone did \textbf{\te{ఈలోగా}};
  \item how it ended and how you felt, with \textbf{\te{చివరికి}} and \textbf{\te{నాకు} … -\te{గా} \te{అనిపించింది}};
  \item one line about what might have been, with \textbf{-\te{ఇ} \te{ఉంటే} … -\te{ఏవాడిని}} or \textbf{-\te{ఏదాన్ని}}.
\end{enumerate}

There is no model on this page. Use only words you already write. The whole exercise is choosing the right past for each line, and joining them.

\begin{tcolorbox}[breakable,colback=teal!4,colframe=teal!35!black,title={Before you move on}]
Now read your lines as you read the village story earlier in this core. Mark every verb \textbf{s} for scene, \textbf{e} for event. Line one should be mostly \textbf{s}; lines two and three mostly \textbf{e}.

Then five spelling checks. A word that starts with long \emph{ī} starts with \textbf{\te{ఈ}}, not \textbf{\te{ఇ}}; one with long \emph{ō} starts with \textbf{\te{ఓ}}, not \textbf{\te{ఒ}}; one with long \emph{ū} starts with \textbf{\te{ఊ}}, not \textbf{\te{ఉ}}. The \emph{au} of \textbf{\te{ఔషధం}} at the start of a word is the letter \textbf{\te{ఔ}}. \textbf{\te{దుఃఖం}} keeps its two rings, \textbf{\te{ః}}.

Last, read line six alone. Does the if-clause end in \textbf{-\te{ఇ} \te{ఉంటే}}, and the main verb in \textbf{-\te{ఏవాడిని}} or \textbf{-\te{ఏదాన్ని}}?
\end{tcolorbox}
\section[(review)]{Review — it depends, and what would have been}

\label{lesson:TE-R325-depends-and-would}



\begin{quote}
Say \textquotedblleft{}to depend\textquotedblright{} (\emph{ādhārapaḍu}) and \textquotedblleft{}that depends on the weather\textquotedblright{} (\emph{adi vātāvaraṇaṁ mīda ādhārapaḍi uṇṭundi}).

\begin{itemize}
\raggedright
  \item \emph{Write it:} \textbf{\te{ః}} and \textbf{\te{ఔ}}, once each
\end{itemize}
\end{quote}

\subsection*{Your turn}
Say these aloud:

\begin{itemize}
\raggedright
  \item Will you come tomorrow? It depends on the rain. If it rains, I will not come.
  \item According to the timetable, the bus leaves at six. If you do not come quickly, it will go.
  \item If I had had time, I would have come.
  \item A woman says: if the bus had come on time, I would have caught the train.
  \item read your six lines aloud to your reader
\end{itemize}

\begin{tcolorbox}[breakable,colback=teal!4,colframe=teal!35!black,title={Before you move on}]
\textquotedblleft{}It depends on …\textquotedblright{}? (\textbf{… \te{మీద} \te{ఆధారపడి} \te{ఉంటుంది}}.) \textquotedblleft{}If … had come\textquotedblright{}? (\textbf{\te{వచ్చి} \te{ఉంటే}}.) \textquotedblleft{}I would have come\textquotedblright{}? (\textbf{\te{వచ్చేవాడిని}}, \textbf{\te{వచ్చేదాన్ని}}.) \textquotedblleft{}According to\textquotedblright{}? (\textbf{\te{ప్రకారం}}.)
\end{tcolorbox}
